\documentclass[11pt]{article}
\usepackage[margin=1in]{geometry}
\usepackage{amsmath,amssymb,amsthm,bm}
\usepackage{booktabs}
\usepackage{enumitem}
\usepackage[hidelinks]{hyperref}

\newtheorem{proposition}{Proposition}
\newtheorem{corollary}{Corollary}
\newtheorem{remark}{Remark}

\newcommand{\R}{\mathbb{R}}
\newcommand{\X}{\bm X}
\newcommand{\Xh}{\widehat{\bm X}}
\newcommand{\W}{\bm W}
\newcommand{\Wh}{\widehat{\bm W}}
\newcommand{\Y}{\bm Y}
\newcommand{\D}{\bm D}
\newcommand{\B}{\bm B}
\newcommand{\Rr}{\bm R}
\newcommand{\Rh}{\widehat{\bm R}}
\newcommand{\Hh}{\widehat{\bm H}}
\newcommand{\Ht}{\widetilde{\bm H}}
\newcommand{\C}{\bm C}
\newcommand{\SigA}{\bm\Sigma_A}
\newcommand{\Qa}{\mathcal Q_A}
\newcommand{\Qw}{\mathcal Q_W}
\newcommand{\tr}{\operatorname{tr}}
\newcommand{\rank}{\operatorname{rank}}
\newcommand{\range}{\operatorname{range}}

\title{\bf Activation-Error Propagation for H-SVDQuant:\\
From Local Reconstruction to Cumulative Block Reconstruction}
\author{Companion note to \texttt{hsvdquant.tex}}
\date{}

\begin{document}
\maketitle

\begin{abstract}
This note studies the sequential calibration objective
\[
  \min\left\|\X_l\W_l-\widehat{\Y}_l\right\|_F^2,
  \qquad
  \widehat{\Y}_l=
  \Qa(\widetilde{\Xh}_l)\Qw(\Rr_l)+\Xh_l\bm L_{1,l}\bm L_{2,l},
\]
where $\X_l$ is the full-precision reference activation and $\Xh_l$ is the
activation produced by already-quantized upstream operators.  The main result
is that the deterministic part of this apparently new problem is exactly the
old H-SVDQuant problem with two substitutions: use the propagated Hessian
$\Hh_l=\Xh_l^\top\Xh_l$, and replace $\W_l$ by a projected cumulative target
$\W_l^\star=\Hh_l^\dagger\Xh_l^\top\X_l\W_l$.  A non-removable projection
error is added to the objective, but it is constant for the current layer.
Consequently the low-rank refit, Hessian deflation, GPTQ feedback, smoothing
gauge, and the joint $F_W+F_A$ code update all retain their form.
\end{abstract}

\section{Reference and propagated paths}

Suppress the block/linear index $l$.  Let
\[
  \X,\Xh\in\R^{m\times c},\qquad \W\in\R^{c\times n},
  \qquad \Y=\X\W.
\]
$\X$ is captured from the untouched full-precision path; $\Xh$ is produced by
the partially quantized path.  Let $\D=\operatorname{diag}(d_1,\ldots,d_c)$ and
$\widetilde{\Xh}=\Xh\D^{-1}$.  To match the objective in the question, write
the full-precision branch in unsmoothed coordinates as
$\B=\bm L_1\bm L_2$ and the quantized residual in smoothed coordinates as
$\Rh=\Qw(\Rr)$.  Then
\begin{equation}\label{eq:deploy}
  \widehat{\Y}
  =\Qa(\widetilde{\Xh})\Rh+\Xh\B.
\end{equation}
The implementation stores the branch in smoothed coordinates,
$\widetilde{\B}=\D\B=\widetilde{\bm L}_1\bm L_2$, so its runtime expression
$(\Xh\D^{-1})\widetilde{\bm L}_1\bm L_2$ is identical to $\Xh\B$.

The old local objective uses $\Xh\W$ as its target.  It therefore teaches the
current layer to preserve its original weights under the new input, but does
not ask it to compensate for upstream error.  The cumulative objective instead
keeps the full-precision target $\X\W$:
\begin{equation}\label{eq:cumulative}
  \mathcal J_{\mathrm{cum}}
  =\left\|\X\W-\Qa(\Xh\D^{-1})\Rh-\Xh\B\right\|_F^2.
\end{equation}

\section{Exact deterministic reduction}

First turn off activation quantization and define the effective unsmoothed
matrix
\[
  \bm A:=\B+\D^{-1}\Rh,
  \qquad \Hh:=\Xh^\top\Xh,
  \qquad \C:=\Xh^\top\X,
  \qquad \bm G:=\X^\top\X.
\]
The deterministic cumulative risk is
\begin{align}
  \mathcal J_0(\bm A)
  &=\|\X\W-\Xh\bm A\|_F^2 \nonumber\\
  &=\tr(\W^\top\bm G\W)-2\tr(\bm A^\top\C\W)
    +\tr(\bm A^\top\Hh\bm A).                         \label{eq:quadratic}
\end{align}

\begin{proposition}[Projected cumulative target]\label{prop:project}
If $\Hh\succ0$, define
\begin{equation}\label{eq:astar}
  \boxed{\quad
  \bm A^\star=\Hh^{-1}\C\W
  =\W+\Hh^{-1}\Xh^\top(\X-\Xh)\W .
  \quad}
\end{equation}
Then
\begin{equation}\label{eq:complete}
  \mathcal J_0(\bm A)
  =\underbrace{\|\Hh^{1/2}(\bm A-\bm A^\star)\|_F^2}_{\text{optimizable by the current layer}}
   +\underbrace{\tr\!\left[\W^\top
     (\bm G-\C^\top\Hh^{-1}\C)\W\right]}_{\mathcal J_{\rm irr}:\ \text{irreducible at this layer}}.
\end{equation}
For singular $\Hh$, the same statement holds with the Moore--Penrose inverse
and $\mathcal J_{\rm irr}=\|(\bm I-\Xh\Xh^\dagger)\X\W\|_F^2$.
\end{proposition}

\begin{proof}
The normal equation of \eqref{eq:quadratic} is
$\Hh\bm A=\C\W$, which gives \eqref{eq:astar}.  Substitute
$\bm A=\bm A^\star+(\bm A-\bm A^\star)$ in \eqref{eq:quadratic}; the cross
term vanishes by the normal equation, yielding \eqref{eq:complete}.  The
singular case is ordinary least squares projected onto $\range(\Xh)$.
\end{proof}

Equation~\eqref{eq:astar} is the cumulative-error corrector.  If
$\Xh=\X+\delta\X$ and the perturbation is small, then to first order
\begin{equation}\label{eq:firstorder}
  \bm A^\star-\W
  \simeq-(\X^\top\X)^{-1}\X^\top\delta\X\,\W.
\end{equation}
Thus the current weight update is the least-squares inverse image of the
upstream activation error.  Error orthogonal to $\range(\Xh)$ cannot be repaired
by any current linear map and belongs to $\mathcal J_{\rm irr}$.

\subsection{Reduction to the original smoothed problem}

Set
\begin{equation}\label{eq:smoothed-target}
  \widetilde{\W}^{\star}:=\D\bm A^\star,
  \qquad \widetilde{\B}:=\D\B,
  \qquad \Ht:=\D^{-1}\Hh\D^{-1}.
\end{equation}
Since $\bm A=\D^{-1}(\widetilde{\B}+\Rh)$,
\begin{equation}\label{eq:standardform}
  \mathcal J_0
  =\mathcal J_{\rm irr}
  +\left\|\Ht^{1/2}
    (\widetilde{\W}^{\star}-\widetilde{\B}-\Rh)\right\|_F^2.
\end{equation}
This is exactly the original H-SVDQuant deterministic objective, with
\begin{equation}\label{eq:substitution}
  \boxed{
  \widetilde{\bm H}=\D^{-1}\X^\top\X\D^{-1}
  \ \longrightarrow\ 
  \Ht=\D^{-1}\Xh^\top\Xh\D^{-1},
  \qquad
  \widetilde{\W}=\D\W
  \ \longrightarrow\ 
  \widetilde{\W}^{\star}=\D\Hh^\dagger\Xh^\top\X\W .}
\end{equation}
The additive constant $\mathcal J_{\rm irr}$ changes reported absolute error
but cannot change the minimizer for the current layer.

\begin{corollary}[Recovery of the local objective]
If $\Xh=\X$, then $\bm A^\star=\W$ and $\mathcal J_{\rm irr}=0$.  All formulas
reduce to the original local H-SVDQuant derivation.
\end{corollary}

\section{Activation quantization and the two-channel risk}

Let
\[
  \bm E_A:=\Qa(\widetilde{\Xh})-\widetilde{\Xh},
  \qquad \mathbb E[\bm E_A]=0,
  \qquad \SigA:=\mathbb E[\bm E_A^\top\bm E_A].
\]
Using \eqref{eq:deploy}, the exact output error is
\[
  \X\W-\widehat{\Y}
  =\left(\X\W-\Xh\bm A^\star\right)
   +\widetilde{\Xh}
    (\widetilde{\W}^{\star}-\widetilde{\B}-\Rh)
   -\bm E_A\Rh.
\]
The first two deterministic terms are orthogonal by the normal equation.
Under the usual zero-mean/independence model for activation quantization noise,
the noise cross terms also vanish.  Therefore
\begin{equation}\label{eq:twochannel}
  \boxed{
  \mathbb E\,\mathcal J_{\mathrm{cum}}
  =\mathcal J_{\rm irr}
   +\underbrace{\tr\!\left(\bm\Delta^\top\Ht\bm\Delta\right)}_{F_W^{\rm cum}}
   +\underbrace{\tr\!\left(\Rh^\top\SigA\Rh\right)}_{F_A^{\rm cum}},
  \qquad
  \bm\Delta=\widetilde{\W}^{\star}-\widetilde{\B}-\Rh .}
\end{equation}

The decomposition keeps its old form, but both channels are now conditioned on
the propagated path: $F_W$ uses $\Ht$ and the shifted target
$\widetilde{\W}^{\star}$, while $F_A$ uses the noise covariance of
$\Qa(\Xh\D^{-1})$.  This matters because compensation can increase the norm or
dynamic range of the residual, reducing deterministic accumulated error while
increasing activation error.  The existing balance coefficient $\lambda$ is
therefore more, not less, important:
\begin{equation}\label{eq:lambda}
  \mathcal J_\lambda
  =F_W^{\rm cum}+\lambda F_A^{\rm cum},
\end{equation}
with $\mathcal J_{\rm irr}$ omitted during optimization because it is constant.

\section{What changes in each H-SVDQuant block}

\subsection{Low-rank refit and Hessian deflation}

Write $\widetilde{\B}=\bm L_1\bm L_2$.  For fixed $\bm L_1$ and $\Rh$, let
$\bm M=\bm L_1^\top\Ht\bm L_1$.  The exact post-code update remains
\begin{equation}\label{eq:l2}
  \bm L_2^\star
  =\bm M^{-1}\bm L_1^\top\Ht
   (\widetilde{\W}^{\star}-\Rh).
\end{equation}
Eliminating $\bm L_2$ gives the same deflated metric
\begin{equation}\label{eq:hperp}
  \Ht_\perp
  =\Ht-\Ht\bm L_1
   (\bm L_1^\top\Ht\bm L_1)^{-1}
   \bm L_1^\top\Ht.
\end{equation}
Hence weighted SVD, the $\Ht$-orthogonal gauge, and rank-$r$ branch updates are
unchanged algebraically.  They are simply applied to
$(\Ht,\widetilde{\W}^{\star})$ rather than
$(\widetilde{\bm H},\widetilde{\W})$.

\subsection{Joint code target}

After eliminating $\bm L_2$, the idealized code block is
\begin{equation}\label{eq:code}
  \min_{\Rh\in\mathcal Q_W}
  \left\|\widetilde{\W}^{\star}-\Rh\right\|_{\Ht_\perp}^2
  +\lambda\|\Rh\|_{\SigA}^2.
\end{equation}
Completing the square gives
\begin{equation}\label{eq:pseudotarget}
  \bm M_Q=\Ht_\perp+\lambda\SigA,
  \qquad
  \boxed{\ \bm T_Q=\bm M_Q^{-1}\Ht_\perp\widetilde{\W}^{\star}\ },
\end{equation}
so GPTQ/GANQ should quantize $\bm T_Q$ under metric $\bm M_Q$.  This is the
same joint $F_W+F_A$ pseudo-target as before, except that the cumulative target
$\widetilde{\W}^{\star}$ replaces the original weight.  Positive activation
noise also regularizes the rank-deficient null space of $\Ht_\perp$.

\subsection{Smoothing}

For fixed $\Xh$ and $\bm A^\star$, changing $\D$ is still a gauge transform:
\[
  \widetilde{\Xh}=\Xh\D^{-1},\qquad
  \widetilde{\W}^{\star}=\D\bm A^\star,\qquad
  \widetilde{\Xh}\widetilde{\W}^{\star}=\Xh\bm A^\star.
\]
Thus $\D$ cannot change the unquantized deterministic optimum or
$\mathcal J_{\rm irr}$.  It still changes weight group ranges and the covariance
$\SigA$ through the quantizer, so the existing $F_W+\lambda F_A$ smoothing
optimization remains valid after replacing the residual relative to $\W$ by
the residual relative to $\bm A^\star$.

\section{Stable finite-sample estimator}

Computing a pseudoinverse from a small cache is unstable.  Let the propagated
streaming Hessian and paired-cache cross-error be
\[
  \widehat{\bm H}=\frac{1}{N}\sum_{t=1}^N\widehat{\bm x}_t
    \widehat{\bm x}_t^\top,
  \qquad
  \widehat{\bm K}=\frac{1}{M}\sum_{t\in\mathcal C}
    \widehat{\bm x}_t(\bm x_t-\widehat{\bm x}_t)^\top.
\]
The implemented estimator is
\begin{equation}\label{eq:estimator}
  \boxed{
  \widehat{\bm A}^{\star}_\varepsilon
  =\W+(\widehat{\bm H}+\varepsilon\bm I)^{-1}
    \widehat{\bm K}\W .}
\end{equation}
It retains the full-token curvature used by GPTQ while estimating only the
upstream mismatch on a bounded, paired reservoir.  The reference and propagated
reservoirs use identical random priorities, so their rows remain aligned.

The damping has a statistical meaning beyond numerical Cholesky stability: it
shrinks compensation in activation directions that are weakly observed under
$\Xh$.  Without it, a small singular value can turn a modest upstream mismatch
into a very large correction, inflate quantization ranges, and increase $F_A$.

\section{Sequential block and linear calibration}

For decoder block $l$, maintain two paths:
\begin{enumerate}[leftmargin=2em]
  \item Run the untouched block on reference input $\X_l$ once.  Cache the
        reference input of every linear and save the full-precision block output
        $\X_{l+1}$.
  \item Run the working block on propagated input $\Xh_l$.  Process true
        dependency groups in order: attention input projections
        $(q,k,v)$, attention output projection, MLP input projections
        (gate, up), and MLP output projection.
  \item Before each group, collect its propagated Hessians/caches using all
        already-quantized predecessor groups.  For each linear, construct
        \eqref{eq:estimator}, run H-SVDQuant, and replace the module immediately.
  \item Run the fully quantized block on $\Xh_l$ to obtain $\Xh_{l+1}$.
\end{enumerate}
Parallel projections are not artificially ordered: $q/k/v$ share the same
input, as do gate/up, so they are calibrated independently from the same cache.
Only genuine graph dependencies are propagated.  This is the block analogue of
GPTQ error feedback, while \eqref{eq:estimator} is its activation-space,
cross-layer counterpart.

\section{Predictions, risks, and ablations}

\begin{table}[h]
\centering
\begin{tabular}{lll}
\toprule
Block input & Linear objective & Interpretation \\
\midrule
reference & local & no inter-block feedback control \\
quantized & local & strategy (a): propagated calibration only \\
reference & cumulative & intra-block accumulated-error compensation \\
quantized & cumulative & strategies (a)+(b), full sequential compensation \\
\bottomrule
\end{tabular}
\end{table}

The theory gives several falsifiable predictions.
\begin{enumerate}[leftmargin=2em]
  \item The cumulative mode should reduce the measured paired target error
        $\|\X\W-\widehat\Y\|_F^2$, but need not reduce the local error
        $\|\Xh\W-\widehat\Y\|_F^2$.
  \item Gains should grow with depth while correction norms remain controlled;
        no gain at the first parallel group is expected because $\Xh_1=\X_1$.
  \item Very large correction-norm ratios diagnose ill-conditioned compensation
        and predict worse weight/activation quantization.  Increasing damping or
        clipping the correction is then preferable to increasing rank.
  \item Because compensation can move mass into the low-precision residual,
        the best $\lambda$ under the cumulative objective may be larger than the
        best $\lambda$ under the local objective.
  \item Report both total paired MSE and its decomposition into upstream MSE,
        projected (irreducible) MSE, quantized state MSE, and correction-norm
        ratio.  Perplexity alone cannot distinguish compensation from overfit.
\end{enumerate}

\section{Summary of theoretical changes}

\begin{center}
\begin{tabular}{p{0.27\linewidth}p{0.30\linewidth}p{0.33\linewidth}}
\toprule
Object & Local reconstruction & Cumulative reconstruction \\
\midrule
Target output & $\Xh\W$ & $\X\W$ \\
Effective weight target & $\W$ & $\bm A^\star=\Hh^\dagger\Xh^\top\X\W$ \\
Curvature & input Hessian of the local cache & $\Hh=\Xh^\top\Xh$ \\
Unavoidable term & zero & $\mathcal J_{\rm irr}$ \\
Smoothed target & $\D\W$ & $\D\bm A^\star$ \\
Low-rank/$\bm L_2$ updates & original formulas & unchanged, with target substitution \\
GPTQ metric & Hessian-deflated metric & unchanged, built from $\Hh$ \\
$F_A$ covariance & noise of $\Qa(\X\D^{-1})$ & noise of $\Qa(\Xh\D^{-1})$ \\
Joint code target & $\bm M_Q^{-1}\bm H_\perp\D\W$ & $\bm M_Q^{-1}\bm H_\perp\D\bm A^\star$ \\
\bottomrule
\end{tabular}
\end{center}

The central message is concise: cumulative activation-error correction does not
require a new quantizer or a new low-rank solver.  It requires a paired reference
path, a propagated path, and the projected target \eqref{eq:astar}.  After that
substitution, the existing H-SVDQuant machinery applies verbatim, with an added
irreducible projection error that must be included when reporting absolute
reconstruction quality.

\end{document}
